\section{Research problem and evidence scope}
\label{sec:problem}

A factor-research harness determines how expression candidates are proposed, evaluated, compared, and retained. Improving that process requires more than increasing the number of formulas produced by a language model. The object of adaptation is the sequence of research decisions, while the object of evaluation is the factor-expression candidate generated under the resulting harness. This study asks whether a language model can improve the former under a fixed protocol, and which sources of information contribute to the measured effect.

Two questions organize the analysis. The first compares a memory-equipped, model-driven process with random legal edits over the same editable surface and with the fixed initial harness. This comparison tests whether model-guided adaptation adds value under matched proposal semantics, decision counts, and candidate evaluation budgets. The second decomposes the model-driven process into scalar feedback, cumulative probe memory, and compact execution diagnostics. It tests whether the additional information channels produce separable improvements over their immediate parent conditions.

All experiments use the development window from 2017-01-01 through 2020-11-30. Task screening and online continuation use probe information only; final select evaluation occurs once after each condition has frozen its champion. The study contains no untouched independent market-period out-of-sample evaluation. Fresh tasks and new development splits within this window provide confirmatory development-panel evidence, and all effect estimates retain that scope.
